% ══════════════════════════════════════════════════════════════════════════════
%  ADVANCED WATER AND RESOURCES — the above-level book (from 2027)
%  Faculty decision, Nov 2026: every student completes the at-level lesson book;
%  students who excel move on to this book. Its topics go deeper into the ideas
%  of the at-level lessons rather than adding more of the same work.
%
%  Each topic: a reading, the deeper idea worked through (I do / We do), then
%  %%PARTD:N%% pulls in the old Part D questions of at-level Lesson N as
%  practice, then a Challenge. Built by make_a5.py (via a5_levels.py).
% ══════════════════════════════════════════════════════════════════════════════

% ══════════════════════════════════════════════════════════════════════════════
%  TOPIC 1 — ENERGY, THE WATER CYCLE AND CLIMATE             (deepens Lesson 1)
% ══════════════════════════════════════════════════════════════════════════════
\setcounter{section}{0}\renewcommand{\lessonart}{1}
\sectiondiv{Energy, the Water Cycle and Climate}{The Sun drives the cycle \textbullet{} Energy in evaporation and condensation \textbullet{} How long water stays \textbullet{} A warmer world}\label{sec:1}
\lessoncodes{Deepens Lesson 1 \textbullet{} VC2S8U09 \textbullet{} VC2S10U11}

\vocabhead
\begin{ctab}
\begin{tabular}{|SW{3.4cm}|W{7.0cm}|}
  \hline\hd & \thead{Word} & \thead{Definition (copy from the slide)} \\ \hline
  & \scaf{\textbf{Store}} & \scaf{A place where water is held for a time: the ocean, ice, groundwater, the air.} \\ \hline
  & \textbf{Flux} & \\ \hline
  & \textbf{Residence time} & \\ \hline
  & \textbf{Transpiration} & \\ \hline
  & \textbf{Groundwater} & \\ \hline
\end{tabular}
\end{ctab}

\lessonpart{A}{Reading and comprehension}{Teal}\label{read:1}

\begin{readingbox}{Reading --- The Sun's Water Pump}
Every year the Sun lifts about half a million cubic kilometres of water from the oceans, lakes, soil and leaves into the air. Evaporation takes energy: each gram of water needs about 2400\,J to escape as vapour, which is why sweating cools you. When the vapour condenses into cloud droplets, that energy is released again, warming the air around it. A thunderstorm is powered by the energy released as its clouds condense.

Scientists describe the water cycle as \textbf{stores}, where water is held, and \textbf{fluxes}, the flows between them. Plants are a big flux: water drawn up through their roots and evaporated from their leaves is \textbf{transpiration}. A single large tree can transpire hundreds of litres on a hot day.

How long water stays in a store is its \textbf{residence time}. A water particle stays in the air for about nine days before it falls as rain. In the deep ocean it may stay for thousands of years. Some \textbf{groundwater}, held in rock deep under inland Australia, fell as rain hundreds of thousands of years ago.

Warmer air can hold more water vapour, about 7\% more for each degree Celsius. As the climate warms, more water evaporates, and when it rains, more can fall at once. Scientists expect both heavier downpours and longer droughts in many places.
\end{readingbox}

\rcreset
\rcq{Where does the energy to evaporate water come from, and where does it go when the vapour condenses?}{3cm}
\rcq{What is the difference between a store and a flux? Give an example of each.}{2.5cm}
\rcq{Why can deep groundwater be called ``fossil water''?}{2.5cm}
\rcq{Why might a warmer climate bring both heavier rain and worse droughts?}{3cm}

\lessonpart{B}{Theory and examples}{Blue}\label{theory:1}

\subsection{Energy in the Water Cycle}\label{sub:a1.1}

\begin{workedex}{A canvas water bag}
Before refrigerators, people hung canvas bags of water on cars and verandahs. Water seeping through the canvas evaporated from the outside, and the water inside stayed cool on a hot day. Explain.

\smallskip
To evaporate, the water particles on the outside of the bag need energy. They take it from the bag and the water inside, so the water inside loses energy and its temperature falls. The hotter and drier the air, the faster the evaporation and the stronger the cooling.
\end{workedex}

\begin{wedo}{Thunderstorms}
Explain, using energy, why a large cloud forming on a summer afternoon can grow into a violent thunderstorm.
\answerbox{4.5cm}
\end{wedo}

\subsection{Residence Time}\label{sub:a1.2}

\begin{skillbox}{Residence time}
\[ \text{residence time} = \frac{\text{amount in the store}}{\text{rate of flow out of the store}} \]
\end{skillbox}

\begin{workedex}{A farm dam}
A farm dam holds 6\,million litres. Cattle, evaporation and leakage remove 40\,000\,L a day, and no rain falls. How long will it last?

\smallskip
$6\,000\,000 \div 40\,000 = \mathbf{150\ days}$: about five months without rain.
\end{workedex}

\begin{wedo}{Water in the air}
The air holds about 13\,000\,km\textsuperscript{3} of water, and about 1400\,km\textsuperscript{3} falls as rain and snow every day. Calculate the residence time of water in the air. Then explain why pumping groundwater that took 100\,000 years to collect is a different problem from using river water.
\workbox{5.5cm}
\end{wedo}

\lessonpart{C}{Practice}{Amber}\label{prac:1}

%%PARTD:1%%

\lessonpart{D}{Challenge}{Violet}\label{ext:1}

\task Most of Melbourne's water comes from rain falling on forested catchments that flows into reservoirs. The climate in these catchments is expected to become warmer and a little drier. Explain, as a chain of causes, how this could reduce the water reaching the reservoirs by more than the drop in rainfall alone. Use the words \emph{evaporation}, \emph{transpiration}, \emph{soil} and \emph{runoff}.
\answerbox{8cm}

\clearpage

% ══════════════════════════════════════════════════════════════════════════════
%  TOPIC 2 — RESOURCES ACROSS THEIR LIFE CYCLE               (deepens Lesson 3)
% ══════════════════════════════════════════════════════════════════════════════
\setcounter{section}{1}\renewcommand{\lessonart}{3}
\sectiondiv{Resources Across Their Life Cycle}{Cradle to grave \textbullet{} Embodied energy \textbullet{} Comparing choices with numbers \textbullet{} The circular economy}\label{sec:2}
\lessoncodes{Deepens Lesson 3 \textbullet{} VC2S8U09 \textbullet{} VC2S10H03}

\vocabhead
\begin{ctab}
\begin{tabular}{|SW{3.4cm}|W{7.0cm}|}
  \hline\hd & \thead{Word} & \thead{Definition (copy from the slide)} \\ \hline
  & \scaf{\textbf{Life cycle}} & \scaf{Every stage of a product's life: getting the raw materials, making it, using it and getting rid of it.} \\ \hline
  & \textbf{Life cycle assessment} & \\ \hline
  & \textbf{Embodied energy} & \\ \hline
  & \textbf{Circular economy} & \\ \hline
  & \textbf{E-waste} & \\ \hline
\end{tabular}
\end{ctab}

\lessonpart{A}{Reading and comprehension}{Teal}\label{read:2}

\begin{readingbox}{Reading --- The Hidden Story of a Can}
An aluminium drink can is in your hand for a few minutes. Its \textbf{life cycle} is far longer. Bauxite ore is mined, often in Queensland or Western Australia, refined into alumina, and then smelted into aluminium using huge amounts of electricity. The metal is rolled, shaped, printed, filled and trucked to a shop. After you drink it, the can is buried, burnt, or recycled.

The total energy used across all these stages is the can's \textbf{embodied energy}. Recycling aluminium saves about 95\% of the energy needed to make new metal from ore, because the hardest step, smelting the ore, is skipped. A \textbf{life cycle assessment} adds up the energy, water, land and pollution of every stage, so that two products, or two choices, can be compared fairly.

In a \textbf{circular economy}, products are designed from the start to be repaired, reused and finally recycled, so materials go round and round instead of ending in landfill. \textbf{E-waste}, old phones, computers and batteries, is a big test: it holds valuable gold, copper and rare metals, but is hard to take apart.
\end{readingbox}

\rcreset
\rcq{List the stages in the life cycle of an aluminium can.}{2.5cm}
\rcq{Why does recycling aluminium save so much energy?}{2.5cm}
\rcq{What does a life cycle assessment add up, and why?}{2.5cm}
\rcq{How is a circular economy different from ``make, use, throw away''?}{2.5cm}

\lessonpart{B}{Theory and examples}{Blue}\label{theory:2}

\subsection{Comparing Choices with Numbers}\label{sub:a2.1}

\begin{workedex}{Glass or plastic?}
A refillable glass bottle takes 5\,MJ (megajoules) to make and 0.2\,MJ to wash each time it is reused. A single-use plastic bottle takes 2\,MJ to make. Compare them over 20 uses.

\smallskip
Glass: $5 + 20 \times 0.2 = 5 + 4 = 9$\,MJ. \quad Plastic: $20 \times 2 = 40$\,MJ.\\
Over 20 uses the glass bottle uses \textbf{less than a quarter} of the energy. Over one use, plastic wins.
\end{workedex}

\begin{wedo}{The keep cup}
A ceramic cup takes 14\,MJ to make and 0.1\,MJ to wash each time. A paper cup takes 0.5\,MJ. After how many uses does the ceramic cup become the better choice? (Hint: find when $14 + 0.1n = 0.5n$.)
\workbox{5cm}
\end{wedo}

\subsection{What a Life Cycle Assessment Leaves Out}\label{sub:a2.2}

\begin{workedex}{Energy is not everything}
A life cycle assessment says paper bags use more energy to make than plastic bags. Does that settle which is better?

\smallskip
No. Energy is only one measure. Plastic bags last for centuries and harm wildlife if they escape into the sea; paper breaks down in months. Paper uses more water and land to grow trees. A fair comparison looks at several impacts, and the answer can depend on how many times each bag is used.
\end{workedex}

\begin{wedo}{Cotton T-shirts}
Growing the cotton for one T-shirt can use about 2700\,L of water. Describe two other stages of a T-shirt's life cycle and an impact of each. Suggest one change that would most reduce its total impact.
\answerbox{5cm}
\end{wedo}

\lessonpart{C}{Practice}{Amber}\label{prac:2}

%%PARTD:3%%

\lessonpart{D}{Challenge}{Violet}\label{ext:2}

\task Australians replace their phones about every two to three years, and most old phones end up in drawers or landfill. Should the law require phone makers to take back every old phone and recover its materials? Make a reasoned case, considering the resources in a phone, the energy and risks of mining them, the cost, and who should pay.
\answerbox{9cm}

\clearpage

% ══════════════════════════════════════════════════════════════════════════════
%  TOPIC 3 — WATER FOOTPRINTS                                (deepens Lesson 4)
% ══════════════════════════════════════════════════════════════════════════════
\setcounter{section}{2}\renewcommand{\lessonart}{4}
\sectiondiv{Water Footprints}{Measuring a household's water \textbullet{} Virtual water in food \textbullet{} Planning to use less}\label{sec:3}
\lessoncodes{Deepens Lesson 4 \textbullet{} VC2S8U09 \textbullet{} VC2S10I04}

\vocabhead
\begin{ctab}
\begin{tabular}{|SW{3.4cm}|W{7.0cm}|}
  \hline\hd & \thead{Word} & \thead{Definition (copy from the slide)} \\ \hline
  & \scaf{\textbf{Water footprint}} & \scaf{The total amount of fresh water used to produce the things a person, family or country uses.} \\ \hline
  & \textbf{Virtual water} & \\ \hline
  & \textbf{Flow rate} & \\ \hline
  & \textbf{Per capita} & \\ \hline
  & \textbf{Target} & \\ \hline
\end{tabular}
\end{ctab}

\lessonpart{A}{Reading and comprehension}{Teal}\label{read:3}

\begin{readingbox}{Reading --- The Water You Eat}
During the Millennium Drought, Melburnians were asked to use no more than 155 litres of water each per day at home. That \textbf{target}, Target 155, counted showers, toilets, washing machines and gardens: the water you can see coming out of a tap.

But far more water is used to make the things we buy. The water used to grow, process and transport a product is its \textbf{virtual water}. It takes roughly 1800 litres to grow a kilogram of wheat, and around 15\,000 litres for a kilogram of beef, mostly as rain and irrigation water for the grass and grain the cattle eat. Adding up the tap water and the virtual water gives a person's \textbf{water footprint}.

Most of a typical Australian's water footprint is in food, not in the shower. So, although shorter showers matter, especially in a drought, what we eat and what we waste make an even bigger difference. Foods grown with rain instead of irrigation, and food that is eaten instead of thrown out, both shrink the footprint.
\end{readingbox}

\rcreset
\rcq{What did Target 155 measure, and what did it leave out?}{2.5cm}
\rcq{What is virtual water? Give an example.}{2.5cm}
\rcq{Why does beef have such a large water footprint?}{2.5cm}
\rcq{Why does wasting food waste water?}{2.5cm}

\lessonpart{B}{Theory and examples}{Blue}\label{theory:3}

\subsection{Measuring a Household's Water}\label{sub:a3.1}

\begin{skillbox}{Water used}
\[ \text{water used (L)} = \text{flow rate (L/min)} \times \text{time (min)} \]
Measure a flow rate by timing how long it takes to fill a bucket of known volume.
\end{skillbox}

\begin{workedex}{A shower}
A shower fills a 9\,L bucket in 60 seconds, so its flow rate is 9\,L/min. An 8-minute shower uses $9 \times 8 = \mathbf{72\ L}$. With a water-saving shower head at 6\,L/min, the same shower uses 48\,L, a saving of 24\,L.
\end{workedex}

\begin{wedo}{One person's day}
Estimate one person's tap water for a day: a 6-minute shower at 9\,L/min; four toilet flushes at 4.5\,L; teeth brushed twice with the tap running for 2 minutes at 6\,L/min; a share of the washing machine, 30\,L. Is it under Target 155?
\workbox{5.5cm}
\end{wedo}

\subsection{Virtual Water in Food}\label{sub:a3.2}

\begin{workedex}{A burger patty}
A 150\,g beef patty, at about 15\,000\,L per kg:\\
$0.15 \times 15\,000 = \mathbf{2250\ L}$, about as much as fifteen days of tap water at Target 155.
\end{workedex}

\begin{wedo}{A sandwich}
Two slices of bread use about 80\,g of wheat (1800\,L per kg). Find the virtual water in the bread, and compare it with the burger patty.
\workbox{4.5cm}
\end{wedo}

\lessonpart{C}{Practice}{Amber}\label{prac:3}

%%PARTD:4%%

\lessonpart{D}{Challenge}{Violet}\label{ext:3}

\task A family of four each use about 200\,L of tap water a day. They want to reach Target 155.
\pt{a} How much less water must the whole family use each day? And each year?
\workbox{4cm}
\pt{b} Plan three changes that together save at least that much. Use numbers for each change.
\answerbox{6cm}

\clearpage

% ══════════════════════════════════════════════════════════════════════════════
%  TOPIC 4 — ENERGY EFFICIENCY AND THE GRID                  (deepens Lesson 5)
% ══════════════════════════════════════════════════════════════════════════════
\setcounter{section}{3}\renewcommand{\lessonart}{5}
\sectiondiv{Energy Efficiency and the Grid}{Efficiency \textbullet{} Kilowatt-hours and cost \textbullet{} Matching supply to demand \textbullet{} Storing energy}\label{sec:4}
\lessoncodes{Deepens Lesson 5 \textbullet{} VC2S10U15 \textbullet{} VC2S10U16}

\vocabhead
\begin{ctab}
\begin{tabular}{|SW{3.4cm}|W{7.0cm}|}
  \hline\hd & \thead{Word} & \thead{Definition (copy from the slide)} \\ \hline
  & \scaf{\textbf{Efficiency}} & \scaf{The percentage of the energy put in that ends up as useful energy out.} \\ \hline
  & \textbf{Power} & \\ \hline
  & \textbf{Kilowatt-hour (kWh)} & \\ \hline
  & \textbf{Demand} & \\ \hline
  & \textbf{Energy storage} & \\ \hline
\end{tabular}
\end{ctab}

\lessonpart{A}{Reading and comprehension}{Teal}\label{read:4}

\begin{readingbox}{Reading --- Keeping the Lights On}
Energy is never destroyed, but it is often wasted. An old-style light globe turns most of its electrical energy into heat, not light. A car engine turns only about a quarter of the energy in its fuel into motion; the rest heats the engine, the exhaust and the air. The \textbf{efficiency} of a device is the share of the energy put in that comes out in the form we want.

\textbf{Power} is how fast energy is used, measured in watts (W) or kilowatts (kW). Electricity companies charge for energy in \textbf{kilowatt-hours} (kWh): one kilowatt of power used for one hour.

The electricity grid has a tricky job: the electricity made must match the electricity used, every second. \textbf{Demand} in Victoria peaks on summer evenings, when air conditioners and ovens are on, just as solar panels stop producing. That is why \textbf{energy storage} matters: batteries, and pumped hydro schemes such as Snowy 2.0, store spare energy from sunny, windy times and release it when it is needed.
\end{readingbox}

\rcreset
\rcq{Where does most of the energy go in an old-style light globe and in a car engine?}{2.5cm}
\rcq{What is the difference between power and energy?}{2.5cm}
\rcq{Why is a summer evening a hard time for a grid with a lot of solar power?}{2.5cm}
\rcq{How does energy storage help?}{2.5cm}

\lessonpart{B}{Theory and examples}{Blue}\label{theory:4}

\subsection{Efficiency}\label{sub:a4.1}

\begin{skillbox}{Efficiency}
\[ \text{efficiency} = \frac{\text{useful energy out}}{\text{total energy in}} \times 100\% \]
No device is 100\% efficient: some energy always ends up as heat, sound or other forms we do not want.
\end{skillbox}

\begin{workedex}{A kettle}
A kettle uses 200\,000\,J of electrical energy, and 180\,000\,J ends up heating the water.

\smallskip
Efficiency $= 180\,000 \div 200\,000 \times 100\% = \mathbf{90\%}$. The other 10\% heats the kettle and the air around it.
\end{workedex}

\begin{wedo}{A car engine}
A car burns fuel containing 400\,kJ of energy, and 100\,kJ becomes motion. Find the efficiency, and name two places the rest of the energy goes.
\workbox{4.5cm}
\end{wedo}

\subsection{Kilowatt-hours and Cost}\label{sub:a4.2}

\begin{skillbox}{Energy and cost}
\[ \text{energy (kWh)} = \text{power (kW)} \times \text{time (h)} \qquad \text{cost} = \text{energy (kWh)} \times \text{price per kWh} \]
1\,kW $=$ 1000\,W.
\end{skillbox}

\begin{workedex}{A heater}
A 2\,kW heater runs for 3 hours. Energy $= 2 \times 3 = 6$\,kWh. At 30 cents per kWh, it costs $6 \times 30 = 180$ cents, or \textbf{\$1.80}.
\end{workedex}

\begin{wedo}{Two light globes}
An old globe uses 60\,W and an LED gives the same light using 9\,W. Each is on 5 hours a day for a year (365 days). Find the energy each uses in kWh, and how much the LED saves at 30 cents per kWh.
\workbox{6cm}
\end{wedo}

\lessonpart{C}{Practice}{Amber}\label{prac:4}

%%PARTD:5%%

\lessonpart{D}{Challenge}{Violet}\label{ext:4}

\task A small town wants all its electricity from rooftop solar. Its solar panels make more than the town needs from 10\,am to 3\,pm, but nothing after sunset, and its demand is highest from 6\,pm to 9\,pm. Explain the problem, propose a solution using energy storage, and describe one advantage and one disadvantage of your solution.
\answerbox{9cm}

\clearpage
